%% 1-12 かけ算は「同じ かずの まとまり」（10ページ、p.120〜129）
%% 第1巻のきまり：15 pt・行送り28 pt、分かち書き（文節の間に半角の空白）、漢字はすべてルビ、式は全角の × ＝ で書く
%% おはじきの ならべ方（1-12・1-13・1-14 で そろえる）：「1つ分」は たての 1れつ、「いくつ分」は れつの かず（よこ に ならぶ）。
%%   こうすると、 3 × 4 は たて 3こ・よこ 4れつ の 四角になり、1-21 の「たて × よこ」とも そろう。

\makeatletter
% おさらと おはじき  \bkqkplate{(x,y)}{おはじきの かず（1〜5）}{色}   (x,y) は おさらの まん中
\newcommand\bkqkplate[3]{\begin{scope}[shift={#1}]
  \pgfmathsetmacro\bkqk@w{max(#2,3)*5.6+9}%
  \bkplate{(0,0)}{\bkqk@w}
  \ifnum#2>0 \foreach \bkqk@k in {1,...,#2} {\bkohajiki{({(\bkqk@k-(#2+1)/2)*5.6},1.2)}{2.4}{#3}}\fi
\end{scope}}
% おはじきの 四角  \bkqkarray{(x,y)}{たての かず}{れつの かず}{間}{半径}{色}   (x,y) は 左下の おはじきの まん中
\newcommand\bkqkarray[6]{\begin{scope}[shift={#1}]
  \foreach \bkqk@i in {1,...,#3} {\foreach \bkqk@j in {1,...,#2} {\bkohajiki{({(\bkqk@i-1)*#4},{(\bkqk@j-1)*#4})}{#5}{#6}}}
\end{scope}}
% たての 1れつ（1つ分）を かこむ わく  \bkqkcols[わくの すきま]{(x,y)}{たての かず}{れつの かず}{間}{半径}
\newcommand\bkqkcols[6][1.2]{\begin{scope}[shift={#2}]
  \foreach \bkqk@i in {1,...,#4} {\draw[bkTheme!75, line width=0.8pt, dash pattern=on 2pt off 1.5pt, rounded corners=1.8mm]
    ({(\bkqk@i-1)*#5-#6-#1},{-#6-#1}) rectangle ({(\bkqk@i-1)*#5+#6+#1},{(#3-1)*#5+#6+#1});}
\end{scope}}
% あめの ふくろ（あめ 5こ）  \bkqkbag{(x,y)}   (x,y) は ふくろの 下の まん中
\newcommand\bkqkbag[1]{\begin{scope}[shift={#1}]
  \filldraw[fill=bkBlue!7, draw=bkBlue!55, line width=0.7pt, rounded corners=1.5mm] (-6.5,0) -- (6.5,0) -- (5.5,13) -- (-5.5,13) -- cycle;
  \fill[bkRed] (-2.6,13.3) -- (0,12.6) -- (2.6,13.3) -- (2.6,14.6) -- (0,13.9) -- (-2.6,14.6) -- cycle;
  \foreach \bkqk@x/\bkqk@y/\bkqk@c in {-2.7/2.6/bkAppleRed,2.7/2.6/bkYellow,-2.7/6.1/bkRoute4,2.7/6.1/bkRoute2,0/9.6/bkRoute5} {\bkcandy{(\bkqk@x,\bkqk@y)}{1.25}{\bkqk@c}}
\end{scope}}
% たまご  \bkqkegg{(x,y)}
\newcommand\bkqkegg[1]{\filldraw[fill=white, draw=bkStem!45, line width=0.5pt] #1 ellipse (2.6 and 3.4);}
% 囲みの 中の 文を、 右で そろえない（字の あいだが 広がらない ように）。 囲みの 中だけで きく
\newcommand\bkqkrag{\rightskip=0pt plus 1fil\relax}
% かんがえかたの 2行めを、 1行めの 書きだしに そろえる 字下げ（「かんがえかた」の 札の はば）
\newcommand\bkqkkindent{\setbox0=\hbox{\tikz\node[font=\sffamily\bfseries\small, inner xsep=5pt, inner ysep=2.5pt]{かんがえかた};}\hspace*{\dimexpr\wd0+0.6\zw\relax}}
% うえから 見た 車  \bkqkcartop{(x,y)}{色}
\newcommand\bkqkcartop[2]{\begin{scope}[shift={#1}]
  \fill[#2, rounded corners=1.2mm] (-2.4,-3.6) rectangle (2.4,3.6);
  \fill[white!75!bkBlue, rounded corners=0.5mm] (-1.7,0.8) rectangle (1.7,2.3);
  \fill[white!75!bkBlue, rounded corners=0.5mm] (-1.7,-2.8) rectangle (1.7,-1.9);
\end{scope}}
% ロボット（よこから）  \bkqkrobot{(x,y)}{色}   (x,y) は 足もとの まん中
\newcommand\bkqkrobot[2]{\begin{scope}[shift={#1}]
  \fill[bkInk!75] (-3.6,1.8) circle (1.8) (3.6,1.8) circle (1.8);
  \fill[#2, rounded corners=1.2mm] (-5,3) rectangle (5,11);
  \fill[#2!70!black, rounded corners=1mm] (-3.6,11.6) rectangle (3.6,17);
  \fill[white] (-1.6,14.4) circle (1) (1.6,14.4) circle (1);
  \fill[bkInk] (-1.4,14.4) circle (0.45) (1.8,14.4) circle (0.45);
  \draw[bkInk!70, line width=0.8pt] (0,17) -- (0,19); \fill[bkRed] (0,19.4) circle (0.8);
  \fill[white, opacity=0.5, rounded corners=0.5mm] (-3,5) rectangle (3,8.5);
\end{scope}}
\makeatother

\bkchapter{1}{12}{かけ\R{算}{ざん}は\\「\R{同}{おな}じ かずの まとまり」}{\bklevel{1}}{かけざんは「おなじ かずの まとまり」}
\bkmichishirube{1-4,1-6}{1-13,1-14,1-21,2-3}

\begin{goal}
\begin{itemize}[label={\color{bkTheme}●}, leftmargin=1.4\zw, itemsep=0pt, topsep=0pt]
\item 「1つ\R{分}{ぶん} × いくつ\R{分}{ぶん}」の しきが かける。
\item かけ\R{算}{ざん}を、 おはじきの \R{四角}{し|かく}で あらわせる。
\item 2のだん〜5のだんの \R{九九}{く|く}を つかえる。
\end{itemize}
\end{goal}

\begin{hatena}
\begin{bkfigure}
\begin{tikzpicture}[x=1mm, y=1mm]
  % A：おさらに 3こずつ
  \fill[white, rounded corners=3mm] (0,0) rectangle (56,46);
  \node[circle, fill=bkYellow, text=white, font=\sffamily\bfseries, minimum size=6.5mm, inner sep=0pt] at (6,40) {A};
  \node[anchor=west, font=\small, text=bkInk] at (11,40) {おさらに 3こずつ};
  \foreach \x/\y in {14.5/25,41.5/25,14.5/9,41.5/9} {\bkqkplate{(\x,\y)}{3}{bkRoute2}}
  % やじるし
  \draw[bk arrow, bkYellow!85!black, line width=2.2pt] (57.5,23) -- (65,23);
  % B：ならべると 四角
  \begin{scope}[shift={(66,0)}]
  \fill[white, rounded corners=3mm] (0,0) rectangle (56,46);
  \node[circle, fill=bkYellow, text=white, font=\sffamily\bfseries, minimum size=6.5mm, inner sep=0pt] at (6,40) {B};
  \node[anchor=west, font=\small, text=bkInk] at (11,40) {ならべると…};
  \draw[bkTheme, line width=1.3pt, dash pattern=on 3pt off 2pt, rounded corners=2mm] (11.5,3.5) rectangle (44.5,33.5);
  \bkqkarray{(16,9)}{3}{4}{8}{3.2}{bkRoute2}
  \end{scope}
\end{tikzpicture}
\end{bkfigure}
\bkqkrag\noindent A の おはじきを ならべなおしたら、 B の \mbox{\R{四角}{し|かく}}に なったよ。 ぜんぶで なんこ かな？\\
\mbox{\R{四角}{し|かく}}に すると、 はやく かぞえられるのかな？
\end{hatena}

\clearpage
%% ---------------- p.121 やってみよう ----------------
\begin{yattemiyou}[おはじきを \R{同}{おな}じ かずずつ おこう]\bkqkrag
\begin{enumerate}[label={\bkcircled[bkGreen]{\arabic*}}, leftmargin=2.2\zw, itemsep=2pt, topsep=0pt]
\item おさらを 4まい ならべて、 どの おさらにも おはじきを 3こずつ のせよう。
\item おさらを ゆびで さしながら、 「3、 6、 9、 …」と かぞえよう。
\end{enumerate}
\begin{bkfigure}
\begin{tikzpicture}[x=1mm, y=1mm]
  \foreach \x in {15,45,75,105} {\bkqkplate{(\x,0)}{3}{bkRoute2}}
  \foreach \x/\t in {15/3,45/6,75/9} {\node[circle, draw=bkGreen, line width=1pt, fill=white, minimum size=8mm, inner sep=0pt, font=\sffamily\bfseries, text=bkGreen!70!black] at (\x,-12) {\t};}
  \draw[bkTheme, line width=1pt, rounded corners=1mm, fill=white] (99,-16.5) rectangle (111,-7.5);
  \foreach \x in {15,45,75} {\draw[bk arrow, bkGreen, line width=1pt] (\x+5.5,-12) -- (\x+23.5,-12);
    \node[font=\sffamily\bfseries\footnotesize, text=bkGreen!70!black] at (\x+14.5,-8.5) {＋3};}
\end{tikzpicture}
\end{bkfigure}
\begin{enumerate}[label={\bkcircled[bkGreen]{\arabic*}}, leftmargin=2.2\zw, itemsep=2pt, topsep=0pt, start=3]
\item 1まいの おさらの 3こを、 たてに 1れつに ならべよう。 4まい\R{分}{ぶん} ならべると、 どんな かたちに なるかな？
\end{enumerate}
\begin{bkfigure}
\begin{tikzpicture}[x=1mm, y=1mm]
  \bkqkcols{(0,0)}{3}{4}{10}{3.4}
  \bkqkarray{(0,0)}{3}{4}{10}{3.4}{bkRoute2}
  \foreach \i in {1,2,3,4} {\node[font=\sffamily\bfseries\footnotesize, text=bkTheme] at ({(\i-1)*10},27.5) {\i};}
  \node[anchor=west, font=\sffamily\bfseries\footnotesize, text=bkTheme] at (35,27.5) {れつ\R{目}{め}};
\end{tikzpicture}
\end{bkfigure}
\noindent ぜんぶで \bkkakikomi\enspace こ。\hspace{2\zw}3こずつの れつが \bkkakikomi\enspace れつ。
\end{yattemiyou}

\bkfukidashi{sora}{3こずつ ふえて いくね！ たてに ならべたら、 \R{四角}{し|かく}に なったよ。}

\clearpage
%% ---------------- p.122〜123 数学の道具 ----------------
\bksection{doguu}{\R{数学}{すう|がく}の \R{道具}{どう|ぐ}}

\bksubsection{1}{\R{同}{おな}じ かずの まとまり}

\label{bk:1-12:matomari}どの おさらにも 3こずつ。 この 3こが 「1つ\R{分}{ぶん}の かず」です。 おさらは 4まい。 3こが 「4つ\R{分}{ぶん}」 あります。

\begin{bkfigure}
\begin{tikzpicture}[x=1mm, y=1mm]
  \foreach \x in {15,45,75,105} {\bkqkplate{(\x,0)}{3}{bkRoute2}}
  \bkbrace{(2,6)}{(28,6)}{1つ\R{分}{ぶん}は 3こ}{1}
  \bkbrace{(2,-5)}{(118,-5)}{4つ\R{分}{ぶん}}{-1}
\end{tikzpicture}
\end{bkfigure}
\noindent ぜんぶの かずは、 たし\R{算}{ざん}で もとめられます。 おさらを 1まいずつ たして いこう。

\begin{bkfigure}
\begin{tikzpicture}[x=1mm, y=1mm]
  \foreach \i/\t/\s in {1/3/{3},2/6/{3 ＋ 3},3/9/{6 ＋ 3},4/12/{9 ＋ 3}} {
    \node[font=\sffamily\bfseries\footnotesize, text=bkInk!70] at ({(\i-1)*30+15},8) {\i まいめ};
    \node[font=\small, text=bkInk] at ({(\i-1)*30+15},0) {\s};
    \node[circle, draw=bkGreen, line width=1pt, fill=white, minimum size=8mm, inner sep=0pt, font=\sffamily\bfseries, text=bkGreen!70!black] at ({(\i-1)*30+15},-9) {\t};}
  \foreach \i in {1,2,3} {\draw[bk arrow, bkGreen, line width=1pt] ({(\i-1)*30+21},-9) -- ({\i*30+9},-9);}
\end{tikzpicture}
\end{bkfigure}
\noindent おさらが 1まい ふえる たびに、 3こ ふえます。 4まいめで 12こ。
\bkbigshiki{3 ＋ 3 ＋ 3 ＋ 3 ＝ 12}{}

\bksubsection{2}{かけ\R{算}{ざん}の しき}

\label{bk:1-12:shiki}3が 4つ\R{分}{ぶん}の ことを、 みじかく こう かきます。\bkidx{ひとつぶん}{1つ\R{分}{ぶん}}\bkidx{いくつぶん}{いくつ\R{分}{ぶん}}\bkidx{かける}{×（かける）}

\noindent\begin{minipage}[c]{0.5\linewidth}\centering
\begin{tikzpicture}[x=1mm, y=1mm]
  \begin{scope}[every node/.style={font=\fontsize{26}{30}\selectfont, inner sep=0pt, anchor=base west}]
  \node[text=bkRoute2] (a) at (-24,-3.2) {3};
  \node (x) at ([xshift=3.2mm]a.base east) {×};
  \node[text=bkOrange] (b) at ([xshift=3.2mm]x.base east) {4};
  \node (e) at ([xshift=3.2mm]b.base east) {＝};
  \node (c) at ([xshift=3.2mm]e.base east) {12};
  \end{scope}
  \draw[bkTheme, line width=1.2pt, rounded corners=3mm] ([xshift=-6mm]a.west |- 0,-7.5) rectangle ([xshift=6mm]c.east |- 0,7.5);
  \draw[bkRoute2, line width=1pt] (a.south |- 0,-8.6) -- ++(-3,-4) coordinate (la);
  \node[anchor=north, font=\sffamily\bfseries\small, text=bkRoute2!85!black, inner sep=1pt] at (la) {1つ\R{分}{ぶん}};
  \draw[bkOrange, line width=1pt] (b.south |- 0,-8.6) -- ++(3,-4) coordinate (lb);
  \node[anchor=north, font=\sffamily\bfseries\small, text=bkOrange!85!black, inner sep=1pt] at (lb) {いくつ\R{分}{ぶん}};
\end{tikzpicture}
\end{minipage}\hfill
\begin{minipage}[c]{0.48\linewidth}\raggedleft
\bkkigocard{×}{かける}{\R{同}{おな}じ かずの いくつ\R{分}{ぶん}}{times}{1-12}
\end{minipage}
\par\vspace{4pt}
\noindent 「3 かける 4 は 12」と よみます。 この けいさんを かけ\R{算}{ざん}と いいます。\bkidx{かけざん}{かけ\R{算}{ざん}}

\noindent ×の \R{左}{ひだり}に 「1つ\R{分}{ぶん}の かず」、 \R{右}{みぎ}に 「いくつ\R{分}{ぶん}」を かきます。 「3こずつ」の ように、 「〜ずつ」が ついた かずが 1つ\R{分}{ぶん}です。 （\R{左}{ひだり}と \R{右}{みぎ}を \R{入}{い}れかえると どう なるかは、 1-13で しらべるよ。）

\bksubsection{3}{ならべると \R{四角}{し|かく}}

\label{bk:1-12:shikaku}\noindent\begin{minipage}[c]{0.58\linewidth}
1さらの 3こを たてに 1れつ。 4さら\R{分}{ぶん}で 4れつ。 おはじきが \mbox{\R{四角}{し|かく}}に ならびます。\bkidx{しかくにならべる}{\R{四角}{し|かく}に ならべる} \mbox{\R{一目}{ひと|め}}で \mbox{3 × 4} と わかるね。\par
\bkhatenaref の おはじきも、 \mbox{3 × 4 ＝ 12}で 12こです。
\end{minipage}\hfill
\begin{minipage}[c]{0.4\linewidth}\centering
\begin{tikzpicture}[x=1mm, y=1mm]
  \bkqkcols{(0,0)}{3}{4}{8.5}{3}
  \bkqkarray{(0,0)}{3}{4}{8.5}{3}{bkRoute2}
  \draw[bkRoute2!80!black, line width=0.7pt, decorate, decoration={brace, amplitude=2mm}] (-5.5,-4.2) -- (-5.5,21.2)
    node[midway, left=3mm, font=\sffamily\bfseries\footnotesize, text=bkRoute2!85!black, align=center] {たて\\3こ};
  \bkbrace{(-4.2,-6)}{(29.7,-6)}{4れつ}{-1}
\end{tikzpicture}
\end{minipage}

\bksubsection{4}{かけ\R{算}{ざん}だと はやく わかる}

\noindent \bkhatenaref の 12こを、 3つの やりかたで かぞえて くらべよう。

\begin{bkfigure}
\begin{tblr}{colspec={Q[l,m,wd=27mm]Q[l,m,wd=78mm]Q[c,m,wd=19mm]}, hlines={0.6pt,bkLine}, vlines={0.6pt,bkLine},
  row{1}={bg=bkTheme!12, font=\sffamily\bfseries\small}, column{1}={font=\sffamily\bfseries\small}, rowsep=2pt, colsep=3pt, cells={font=\small}}
かぞえかた & いう かず & かいすう \\
① 1こずつ & 1、 2、 3、 4、 5、 6、 7、 8、 9、 10、 11、 12 & 12かい \\
② 3こずつ & 3、 6、 9、 12 & 4かい \\
③ かけ\R{算}{ざん} & \mbox{3 × 4 ＝ 12} & 1かい \\
\end{tblr}
\end{bkfigure}
\noindent ①は 12かい、 ②は 4かい、 ③は 1かいで わかります。 \R{四角}{し|かく}に ならべると、 「3こが 4れつ」と \R{一目}{ひと|め}で わかるので、 かけ\R{算}{ざん}が すぐ つかえるのです。

\noindent おはじきが もっと おおいと、 ちがいは もっと \R{大}{おお}きく なります。 3こずつ 9さらなら、 たし\R{算}{ざん}では 3を 9こも ならべる ことに なります。
\par\vspace{2pt}
\noindent\hfill{\small 3 ＋ 3 ＋ 3 ＋ 3 ＋ 3 ＋ 3 ＋ 3 ＋ 3 ＋ 3 ＝ 27}\hfill\null
\par\vspace{2pt}
\noindent かけ\R{算}{ざん}なら、 \mbox{3 × 9 ＝ 27} の 1つの しきで すみます。

\bksubsection{5}{いくつ\R{分}{ぶん}が 1 ふえると}

\label{bk:1-12:fueru}\noindent \mbox{3 × 1} は、 おさらが 1まい だけ。 だから こたえは 3です。 おさらが 1まい ふえると、 3こ ふえます。

\begin{bkfigure}
\begin{tikzpicture}[x=1mm, y=1mm]
  \bkqkcols{(4,0)}{3}{1}{8}{3}
  \bkqkarray{(4,0)}{3}{1}{8}{3}{bkRoute2}
  \node[font=\sffamily\bfseries\footnotesize] at (4,-9) {3 × 1 ＝ 3};
  \bkqkcols{(36,0)}{3}{4}{8}{3}
  \bkqkarray{(36,0)}{3}{4}{8}{3}{bkRoute2}
  \draw[bkOrange, line width=1pt, dash pattern=on 2pt off 1.5pt, rounded corners=1.8mm] (64.8,-4.2) rectangle (75.2,20.2);
  \bkqkarray{(70,0)}{3}{1}{8}{3}{bkOrange}
  \node[font=\sffamily\bfseries\footnotesize] at (48,-9) {3 × 4 ＝ 12};
  \node[font=\sffamily\bfseries\footnotesize, text=bkOrange!85!black] at (70,-9) {＋ 3};
  \draw[bk arrow, bkTheme, line width=1.2pt] (82,8) -- (92,8);
  \node[anchor=west, font=\sffamily\bfseries\small, text=bkTheme] at (94,8) {3 × 5 ＝ 15};
  \node[anchor=west, font=\sffamily\footnotesize, text=bkInk] at (94,0.5) {12 ＋ 3 ＝ 15};
\end{tikzpicture}
\end{bkfigure}
\noindent \mbox{3 × 5} は、 \mbox{3 × 4} より 1れつ（3こ） おおい。 だから \mbox{3 × 5 ＝ 12 ＋ 3 ＝ 15} です。 こたえを わすれても、 1つ \R{前}{まえ}の こたえに 1つ\R{分}{ぶん}を たせば もとめられます。

\bksubsection{6}{2のだん〜5のだん}

\label{bk:1-12:dan}\mbox{3 × 1}、 \mbox{3 × 2}、 \mbox{3 × 3}、 … と ならべた ものを、 3のだんの \mbox{\R{九九}{く|く}}と いいます。\bkidx{くく}{\R{九九}{く|く}} 1れつ ふえると、 こたえは 3 ふえます。

\begin{bkfigure}
\begin{tikzpicture}[x=1mm, y=1mm]
  \foreach \n/\c in {1/13,2/45,3/79,4/114} {
    \bkqkarray{({\c-(\n-1)*2.75},0)}{3}{\n}{5.5}{2.1}{bkRoute2}
    \pgfmathtruncatemacro\a{3*\n}
    \node[font=\sffamily\bfseries\footnotesize, text=bkInk] at (\c,-6.5) {3 × \n ＝ \a};}
  \foreach \a/\b in {17.1/38.2,51.9/69.4,88.6/101.7} {\draw[bk arrow, bkTheme, line width=1pt] (\a,5.5) -- (\b,5.5);
    \node[font=\sffamily\bfseries\footnotesize, text=bkTheme] at ({(\a+\b)/2},9.5) {＋3};}
\end{tikzpicture}
\end{bkfigure}

\begin{bkfigure}
\begin{tblr}{colspec={Q[c,m,wd=20mm]*{9}{Q[c,m,wd=10.5mm]}}, hlines={0.6pt,bkLine}, vlines={0.6pt,bkLine},
  row{1}={bg=bkTheme!12, font=\sffamily\bfseries\small}, column{1}={font=\sffamily\bfseries\small}, cell{1}{1}={bg=white},
  rowsep=1pt, colsep=0pt, cells={font=\small}}
 & × 1 & × 2 & × 3 & × 4 & × 5 & × 6 & × 7 & × 8 & × 9 \\
2のだん & 2 & 4 & 6 & 8 & 10 & 12 & 14 & 16 & 18 \\
3のだん & 3 & 6 & 9 & 12 & 15 & 18 & 21 & 24 & 27 \\
4のだん & 4 & 8 & 12 & 16 & 20 & 24 & 28 & 32 & 36 \\
5のだん & 5 & 10 & 15 & 20 & 25 & 30 & 35 & 40 & 45 \\
\end{tblr}
\end{bkfigure}

\noindent \R{表}{ひょう}を よこに \R{見}{み}よう。 2のだんは 2ずつ、 4のだんは 4ずつ、 5のだんは 5ずつ ふえて います。 どの だんも、 \mbox{「1つ\R{分}{ぶん}の かず」ずつ} ふえるのです。 だから、 \R{表}{ひょう}は ⑤の 「1つ \R{前}{まえ}の こたえに 1つ\R{分}{ぶん}を たす」で、 \R{自分}{じ|ぶん}で つくれます。

\begin{taisetsu}
かけ\R{算}{ざん}は 「1つ\R{分}{ぶん}の かず × いくつ\R{分}{ぶん}」。\\
おはじきを ならべると、 \R{四角}{し|かく}に なります。
\end{taisetsu}

%% ---------------- p.124〜125 例題 ----------------

\begin{reidai}[1つ\R{分}{ぶん}と いくつ\R{分}{ぶん}]
あめが 1ふくろに 5こずつ。 3ふくろで なんこ？

\noindent\begin{minipage}[c]{0.7\linewidth}
\kangaekata 1つ\R{分}{ぶん}の 5こが、 3つ\R{分}{ぶん}。

\shiki 5 × 3 ＝ 15 \kotaenext 15こ
\end{minipage}\hfill
\begin{minipage}[c]{0.28\linewidth}\centering
\begin{tikzpicture}[x=0.82mm, y=0.82mm]
  \foreach \x in {0,13.5,27} {\bkqkbag{(\x,0)}}
\end{tikzpicture}
\end{minipage}
\end{reidai}

\begin{jibun}
ケーキが 1はこに 2こずつ。 4はこで なんこ？
\bkshikiran
\bkjans{\mbox{2 × 4 ＝ 8}\quad 8こ}
\end{jibun}

\begin{reidai}[\R{四角}{し|かく}に ならんだ おはじき]
\noindent\begin{minipage}[c]{0.7\linewidth}
おはじきは ぜんぶで なんこ？

\kangaekata たて 4こが 6れつ。

\shiki 4 × 6 ＝ 24 \kotaenext 24こ
\end{minipage}\hfill
\begin{minipage}[c]{0.28\linewidth}\centering
\begin{tikzpicture}[x=0.85mm, y=0.85mm]
  \bkqkcols[0.8]{(0,0)}{4}{6}{6.8}{2.1}
  \bkqkarray{(0,0)}{4}{6}{6.8}{2.1}{bkRoute2}
\end{tikzpicture}
\end{minipage}
\end{reidai}

\begin{jibun}
\noindent\begin{minipage}[c]{0.66\linewidth}
おはじきは ぜんぶで なんこ？
\end{minipage}\hfill
\begin{minipage}[c]{0.32\linewidth}\centering
\begin{tikzpicture}[x=1mm, y=1mm]
  \bkqkarray{(0,0)}{3}{7}{5.6}{2.2}{bkRoute4}
\end{tikzpicture}
\end{minipage}
\bkshikiran
\bkjans{\mbox{3 × 7 ＝ 21}\quad 21こ}
\end{jibun}

\begin{reidai}[たし\R{算}{ざん}を かけ\R{算}{ざん}に]
\mbox{3 ＋ 3 ＋ 3 ＋ 3 ＋ 3} を、 かけ\R{算}{ざん}の しきに しよう。

\par\vspace{1pt}\noindent\hfill
\begin{tikzpicture}[x=1mm, y=1mm]
  \bkruler{(0,0)}{7.4}{15}
  \foreach \s in {0,3,6,9,12} {\draw[bkBlue, line width=1.3pt, -{Stealth[length=5pt]}] (\s*7.4+0.6,4) .. controls (\s*7.4+5,10.5) and (\s*7.4+17.2,10.5) .. (\s*7.4+21.6,4.2);}
  \fill[bkBlue] (0,3.2) circle (1.6);
\end{tikzpicture}\hfill\null\par\vspace{1pt}
\kangaekata 3が 5つ\R{分}{ぶん}。 3ずつ 5かい すすむと 15。

\shiki 3 × 5 ＝ 15 \kotaenext 15
\end{reidai}

\begin{jibun}
\mbox{4 ＋ 4 ＋ 4} を、 かけ\R{算}{ざん}の しきに しよう。
\bkshikiran[3\zw]
\bkjans{\mbox{4 × 3 ＝ 12}\quad 12}
\end{jibun}

\begin{reidai}[\R{長}{なが}さの かけ\R{算}{ざん}]
4\,cmの テープを 5\R{本}{ほん} つなぐと、 なん\,cm？
\par\vspace{1pt}\noindent\hfill
\begin{tikzpicture}[x=1mm, y=1mm]
  \foreach \i in {0,...,4} {\bktape{(\i*17,0)}{17}{bkRoute7}{4\,cm}}
  \node[anchor=west, font=\sffamily\bfseries\small, text=bkInk] at (87,3.5) {＝ なん\,cm？};
\end{tikzpicture}\hfill\null\par\vspace{1pt}
\kangaekata 1つ\R{分}{ぶん}は 4\,cm、 いくつ\R{分}{ぶん}は 5\R{本}{ほん}。

\shiki 4 × 5 ＝ 20 \kotaenext 20\,cm
\end{reidai}

\begin{jibun}
3\,cmの テープを 6\R{本}{ほん} つなぐと、 なん\,cm？
\bkshikiran
\bkjans{\mbox{3 × 6 ＝ 18}\quad 18\,cm}
\end{jibun}

\begin{reidai}[こたえを わすれたら]
\mbox{4 × 7} の こたえを わすれました。 \mbox{4 × 6 ＝ 24} は おぼえて います。 \mbox{4 × 7} は いくつ？
\par\vspace{2pt}\noindent\hfill
\begin{tikzpicture}[x=1mm, y=1mm]
  \bkqkcols[0.8]{(0,0)}{4}{6}{7}{2.4}
  \bkqkarray{(0,0)}{4}{6}{7}{2.4}{bkRoute2}
  \draw[bkOrange, line width=0.8pt, dash pattern=on 2pt off 1.5pt, rounded corners=1.6mm] (45.4,-3.2) rectangle (52.6,24.2);
  \bkqkarray{(49,0)}{4}{1}{7}{2.4}{bkOrange}
  \node[anchor=west, font=\sffamily\bfseries\footnotesize, text=bkInk] at (60,14) {4 × 6 ＝ 24};
  \node[anchor=west, font=\sffamily\bfseries\footnotesize, text=bkOrange!85!black] at (60,6) {＋ 4（1れつ ふえた）};
\end{tikzpicture}\hfill\null\par\vspace{2pt}
\kangaekata 1れつ（4こ） ふえる。\\
\bkqkkindent\mbox{24 ＋ 4 ＝ 28}

\shiki 4 × 7 ＝ 28 \kotaenext 28
\end{reidai}

\begin{jibun}
\mbox{5 × 7 ＝ 35} です。 \mbox{5 × 8} は いくつ？ \mbox{5 × 7} の こたえに たして もとめよう。
\bkshikiran[3\zw]
\bkjans{\mbox{35 ＋ 5 ＝ 40}\quad \mbox{5 × 8 ＝ 40}}
\end{jibun}

\begin{reidai}[しきから もんだいを つくろう]
\mbox{2 × 6} の しきに なる もんだいを つくろう。

\kangaekata 2が 1つ\R{分}{ぶん}、 6が いくつ\R{分}{ぶん}。 「2こずつ」が 「6つ」 ある ばめんを かんがえる。

\kotae （れい）1さらに いちごが 2こずつ。 6さらで なんこ？
\end{reidai}

\begin{jibun}
\mbox{5 × 3} の しきに なる もんだいを つくろう。 （ちがう もんだいでも、 「5こずつ」が 「3つ」 なら せいかい。）
\bkwrite{2}
\bkjans{（れい）あめが 1ふくろに 5こずつ。 3ふくろで なんこ？\quad こたえは 15こ}
\end{jibun}

\begin{reidai}[どちらが なんこ おおい？]
A は 2こずつ 6さら、 B は 3こずつ 3さら。 どちらが なんこ おおい？
\par\vspace{2pt}\noindent\hfill
\begin{tikzpicture}[x=1mm, y=1mm]
  \node[circle, fill=bkYellow, text=white, font=\sffamily\bfseries\footnotesize, minimum size=5mm, inner sep=0pt] at (-7,13) {A};
  \bkqkcols[0.8]{(0,0)}{2}{6}{7}{2.4}
  \bkqkarray{(0,0)}{2}{6}{7}{2.4}{bkRoute2}
  \node[circle, fill=bkYellow, text=white, font=\sffamily\bfseries\footnotesize, minimum size=5mm, inner sep=0pt] at (60,13) {B};
  \bkqkcols[0.8]{(67,0)}{3}{3}{7}{2.4}
  \bkqkarray{(67,0)}{3}{3}{7}{2.4}{bkRoute4}
\end{tikzpicture}\hfill\null\par\vspace{2pt}
\kangaekata ① かけ\R{算}{ざん}で、 それぞれの かずを だす。\\
\bkqkkindent A \mbox{2 × 6 ＝ 12}\quad B \mbox{3 × 3 ＝ 9}\\
\bkqkkindent② ちがいは ひき\R{算}{ざん}（1-5）。 \mbox{12 － 9 ＝ 3}

\kotae A の ほうが 3こ おおい
\end{reidai}

\bkfukidashi{mina}{おさらの かずが おおい ほうが おおいとは かぎらないね。 まず かけ\R{算}{ざん}で それぞれの かずを だしてから くらべよう。}

\begin{jibun}
A は 5こずつ 2はこ、 B は 2こずつ 4はこ。 どちらが なんこ おおい？
\bkshikiran[3\zw]
\bkshikiran[3\zw]
\par\vspace{2pt}\noindent\hfill\bkkakikomi\enspace の ほうが \bkkakikomi\enspace こ おおい。
\bkjans{\mbox{5 × 2 ＝ 10}、 \mbox{2 × 4 ＝ 8}、 \mbox{10 － 8 ＝ 2}\quad A の ほうが 2こ おおい}
\end{jibun}

%% ---------------- p.126 まちがえやすいところ ----------------
\bksection{tsumazuki}{まちがえやすい ところ}

\begin{tsumazuki}{かずを たして しまう}\bkqkrag
1さらに 3こずつ、 4さら。 ぜんぶで なんこ？\\
\bkbatsu\enspace \mbox{3 ＋ 4 ＝ 7}\hspace{3\zw}\bkmaru\enspace \mbox{3 × 4 ＝ 12}\par
1さらに 3こずつ、 それが 4さら\R{分}{ぶん}。 ぜんぶの かずを きいて いるから、 かけ\R{算}{ざん}です。
\end{tsumazuki}

\begin{tsumazuki}{\R{同}{おな}じ かずずつ では ない}\bkqkrag
\noindent\begin{minipage}[c]{0.4\linewidth}
\bkbatsu\enspace \mbox{3 × 3 ＝ 9}\par
\bkmaru\enspace \mbox{3 ＋ 1 ＋ 4 ＝ 8}
\end{minipage}\hfill
\begin{minipage}[c]{0.58\linewidth}\centering
\begin{tikzpicture}[x=0.85mm, y=0.85mm]
  \bkqkplate{(0,0)}{3}{bkRoute2}\bkqkplate{(29,0)}{1}{bkRoute2}\bkqkplate{(59,0)}{4}{bkRoute2}
  \foreach \x/\t in {0/3こ,29/1こ,59/4こ} {\node[font=\sffamily\bfseries\footnotesize, text=bkRed!80!black] at (\x,-6.5) {\t};}
\end{tikzpicture}
\end{minipage}
\par かずが ちがう ときは、 たし\R{算}{ざん}で もとめよう。
\end{tsumazuki}

\begin{tsumazuki}{なんかい たしたか わからない}\bkqkrag
\mbox{4 × 6} は 4を 6かい。 ゆびを おりながら かぞえよう。
\begin{bkfigure}
\begin{tikzpicture}[x=1mm, y=1mm]
  \foreach \i/\t in {1/4,2/8,3/12,4/16,5/20,6/24} {
    \node[font=\sffamily\bfseries, text=bkInk] at (\i*17,0) {\t};
    \node[circle, fill=bkRed!80, text=white, font=\sffamily\bfseries\footnotesize, minimum size=5.5mm, inner sep=0pt] at (\i*17,-7) {\i};}
  \node[anchor=east, font=\sffamily\footnotesize, text=bkRed!80!black] at (10,-7) {ゆび};
\end{tikzpicture}
\end{bkfigure}
\end{tsumazuki}

\begin{tsumazuki}{× 1 を ＋ 1 と まちがえる}\bkqkrag
\bkbatsu\enspace \mbox{3 × 1 ＝ 4}\hspace{3\zw}\bkmaru\enspace \mbox{3 × 1 ＝ 3}\par
\mbox{3 × 1} は 「3こが 1さら\R{分}{ぶん}」。 おさらが 1まい だけ だから、 3こ です。 1を たすのでは ないよ。
\end{tsumazuki}

\begin{column}{「×」の しるし}\bkqkrag
\noindent\begin{minipage}[c]{0.8\linewidth}
「×」は、 400\R{年}{ねん}ぐらい \R{前}{まえ}に イギリスで うまれた しるしです。 コンピュータでは、 「×」の かわりに 「＊」を つかいます。
\end{minipage}\hfill
\begin{minipage}[c]{0.18\linewidth}\centering
\begin{tikzpicture}[x=1mm, y=1mm]
  \fill[bkGray!35, rounded corners=2mm] (-9,-9.5) rectangle (9,8.5);
  \filldraw[fill=white, draw=bkGray!80, line width=0.8pt, rounded corners=1.8mm] (-7.5,-7) rectangle (7.5,7);
  \node[font=\sffamily\bfseries\fontsize{20}{20}\selectfont, text=bkInk] at (0,-0.8) {＊};
\end{tikzpicture}
\end{minipage}
\end{column}

\begin{column}{\R{九九}{く|く}の となえかた}\bkqkrag
\R{九九}{く|く}は、 みじかい ことばで となえて おぼえます。 2のだんは こうです。
\par\vspace{3pt}
\noindent\hfill\begin{tblr}{colspec={*{3}{Q[c,m,wd=36mm]}}, colsep=2pt, rowsep=0pt, cells={font=\small}, hlines={0.5pt,bkLine}, vlines={0.5pt,bkLine}}
{\mbox{2 × 1 ＝ 2}\\{\color{bkTheme!80!black}にいちが に}} & {\mbox{2 × 2 ＝ 4}\\{\color{bkTheme!80!black}ににんが し}} & {\mbox{2 × 3 ＝ 6}\\{\color{bkTheme!80!black}にさんが ろく}} \\
{\mbox{2 × 4 ＝ 8}\\{\color{bkTheme!80!black}にしが はち}} & {\mbox{2 × 5 ＝ 10}\\{\color{bkTheme!80!black}にご じゅう}} & {\mbox{2 × 6 ＝ 12}\\{\color{bkTheme!80!black}にろく じゅうに}} \\
{\mbox{2 × 7 ＝ 14}\\{\color{bkTheme!80!black}にしち じゅうし}} & {\mbox{2 × 8 ＝ 16}\\{\color{bkTheme!80!black}にはち じゅうろく}} & {\mbox{2 × 9 ＝ 18}\\{\color{bkTheme!80!black}にく じゅうはち}} \\
\end{tblr}\hfill\null
\par\vspace{3pt}
となえかたを わすれても だいじょうぶ。 ⑤の やりかたで、 いつでも つくりなおせるよ。
\end{column}

\begin{yattemiyou}[うちの なかで さがそう]\bkqkrag
うちの なかで、 「\R{同}{おな}じ かずの まとまり」を さがして、 \R{表}{ひょう}に かこう（たまごの パック、 いすの あし など）。
\par\vspace{3pt}
\noindent\hfill\begin{tblr}{colspec={Q[l,m,wd=34mm]Q[c,m,wd=22mm]Q[c,m,wd=22mm]Q[c,m,wd=34mm]}, hlines={0.6pt,bkLine}, vlines={0.6pt,bkLine},
  row{1}={bg=bkGreen!12, font=\sffamily\bfseries\small}, rowsep=3pt, colsep=2pt, cells={font=\small}}
みつけた もの & 1つ\R{分}{ぶん} & いくつ\R{分}{ぶん} & しき \\
くつした & 2こ & 3くみ & \mbox{2 × 3 ＝ 6} \\
\bkkakikomi[30mm] & \bkkakikomi[14mm] & \bkkakikomi[14mm] & \bkkakikomi[30mm] \\
\end{tblr}\hfill\null
\end{yattemiyou}

%% ---------------- p.127 つかわれているところ ----------------
\begin{butsuri}\bkqkrag
\begin{bkfigure}
\begin{tikzpicture}[x=1mm, y=1mm]
  % たまごの パック
  \fill[white, rounded corners=3mm] (0,0) rectangle (39,40);
  \fill[bkGray!18, rounded corners=2mm] (2.5,15) rectangle (36.5,36);
  \foreach \i in {0,...,4} {\foreach \j in {0,1} {\bkqkegg{(7+\i*6.25,20.5+\j*10)}}}
  \node[font=\sffamily\bfseries\footnotesize] at (19.5,9.5) {たまごの パック};
  \node[font=\sffamily\bfseries\footnotesize, text=bkPurple!80!black] at (19.5,4) {2 × 5 ＝ 10};
  % ちゅうしゃじょう
  \begin{scope}[shift={(41.5,0)}]
  \fill[white, rounded corners=3mm] (0,0) rectangle (39,40);
  \fill[bkGray!55, rounded corners=1.5mm] (3,13) rectangle (36,38);
  \foreach \i in {0,...,5} {\draw[white, line width=0.8pt] (4.5+\i*6,14) -- (4.5+\i*6,37);}
  \foreach \j in {1,2} {\draw[white, line width=0.8pt] (4.5,14+\j*7.7) -- (34.5,14+\j*7.7);}
  \foreach \i/\c in {0/bkRed,1/bkBlue,2/bkYellow,3/bkGreen,4/bkPurple} {\foreach \j in {0,1,2} {\bkqkcartop{(7.5+\i*6,17.85+\j*7.7)}{\c}}}
  \node[font=\sffamily\bfseries\footnotesize] at (19.5,9.5) {ちゅうしゃじょう};
  \node[font=\sffamily\bfseries\footnotesize, text=bkPurple!80!black] at (19.5,4) {3 × 5 ＝ 15};
  \end{scope}
  % ロボット
  \begin{scope}[shift={(83,0)}]
  \fill[white, rounded corners=3mm] (0,0) rectangle (39,40);
  \draw[bkInk!60, line width=0.8pt] (2,15) -- (37,15);
  \bkqkrobot{(7,15)}{bkRoute6}
  \foreach \i in {0,...,3} {\draw[-{Stealth[length=3pt]}, bkTheme, line width=0.9pt] (14+\i*5.5,19) -- (19.3+\i*5.5,19);}
  \foreach \i in {0,...,4} {\draw[bkInk!60, line width=0.5pt] (14+\i*5.5,17) -- (14+\i*5.5,21);}
  \node[font=\sffamily\bfseries\fontsize{8}{9}\selectfont, text=bkTheme] at (25,24.5) {5\,cmを 4かい};
  \node[font=\sffamily\bfseries\footnotesize] at (19.5,9.5) {ロボット};
  \node[font=\sffamily\bfseries\footnotesize, text=bkPurple!80!black] at (19.5,4) {5 × 4 ＝ 20};
  \end{scope}
\end{tikzpicture}
\end{bkfigure}
\begin{itemize}[label={\color{bkPurple}●}, leftmargin=1.4\zw, itemsep=0pt, topsep=0pt]
\item たまごは 2こずつ 5れつで 10こ。
\item ちゅうしゃじょうは 3だいずつ 5れつで 15だい。
\item \R{工場}{こう|じょう}の ロボットは \R{同}{おな}じ うごきを くりかえす。 5\,cmずつ 4かい すすむと \mbox{5 × 4 ＝ 20} で 20\,cm。
\end{itemize}
\par\vspace{2pt}
\noindent\begin{minipage}[c]{0.3\linewidth}\centering
\begin{tikzpicture}[x=1mm, y=1mm]
  \bkclock{(0,0)}{15}{9}{20}
\end{tikzpicture}
\end{minipage}\hfill
\begin{minipage}[c]{0.68\linewidth}
\R{時計}{と|けい}の \R{長}{なが}い はりは、 \R{数字}{すう|じ} 1つで 5\R{分}{ふん}（1-10）。 \R{長}{なが}い はりが 4を さして いたら、 5\R{分}{ふん}が 4つ\R{分}{ぶん}で \mbox{5 × 4 ＝ 20}。 20\R{分}{ぷん}です。 \R{時計}{と|けい}にも 5のだんが かくれて いるよ。
\end{minipage}
\end{butsuri}

\bkfukidashi{delta}{「1かいで すすむ \R{長}{なが}さ × かいすう」は、 この さきの \R{本}{ほん}で 「はやさ × \R{時間}{じ|かん}」に なるよ。 \R{車}{くるま}や ロケットの うごきも わかるんだ！}

\begin{bkfigure}
\begin{tikzpicture}[x=1mm, y=1mm]
  \draw[bkRoute2, line width=3pt] (10,0) -- (128,0);
  \node[anchor=west, inner sep=0pt] at (0,12) {\bkrouteicon[1.4em]{2}\hspace{0.3em}{\sffamily\bfseries\small\color{bkRoute2!80!black}たしあわせの \R{路線}{ろ|せん}}};
  \foreach \x/\c/\t/\s in {10/bkVol1/{1-4}/{たし\R{算}{ざん}},39.5/bkVol1/{1-7}/{\R{大}{おお}きな かず},69/bkVol1/{1-12}/{かけ\R{算}{ざん}},98.5/bkVol1/{1-21}/{ひろさ},128/bkVol2/{2-3}/{{\ltjsetparameter{jacharrange={+200}}＋－×÷}}} {
    \fill[white, draw=\c, line width=1.6pt] (\x,0) circle (2.6);
    \node[anchor=base, font=\sffamily\bfseries\footnotesize, text=\c!85!black] at (\x,-8.4) {\t};
    \node[anchor=base, font=\sffamily\footnotesize, text=bkInk] at (\x,-15.4) {\s};}
  \node[fill=bkTheme, text=white, rounded corners=1mm, font=\sffamily\bfseries\fontsize{8}{9}\selectfont, inner sep=1.5pt] at (69,5.5) {いま ここ};
\end{tikzpicture}
\end{bkfigure}



\clearpage
%% ---------------- p.128〜129 もんだい ----------------
\bkendproblems

\bkproblemgroup{きほん}{1}

\bkproblem けいさんを しましょう。
\par\vspace{2pt}
\noindent\hspace*{1\zw}\begin{tblr}{colspec={*{2}{Q[l,m,wd=62mm]}}, colsep=0pt, rowsep=2pt}
（1）2 × 8 ＝ \bkkakikomi & （2）3 × 9 ＝ \bkkakikomi \\
（3）4 × 7 ＝ \bkkakikomi & （4）5 × 6 ＝ \bkkakikomi \\
（5）2 × 5 ＝ \bkkakikomi & （6）5 × 9 ＝ \bkkakikomi \\
\end{tblr}
\bkpans{（1）16\quad（2）27\quad（3）28\quad（4）30\quad（5）10\quad（6）45}

\bkproblem えんぴつが 1はこに 5\R{本}{ほん}ずつ。 7はこで なん\R{本}{ぼん}？
\bkshikiran
\bkpans{\mbox{5 × 7 ＝ 35}\quad 35\R{本}{ほん}}

\bkproblem おはじきは ぜんぶで なんこ？\par
\noindent\begin{minipage}[c]{0.3\linewidth}\centering
\begin{tikzpicture}[x=1mm, y=1mm]
  \bkqkarray{(0,0)}{5}{4}{6}{2.3}{bkRoute5}
\end{tikzpicture}
\end{minipage}\hfill
\begin{minipage}[c]{0.68\linewidth}
\bkshikiran
\end{minipage}
\bkpans{\mbox{5 × 4 ＝ 20}\quad 20こ}

\bkproblem \mbox{2 ＋ 2 ＋ 2 ＋ 2 ＋ 2 ＋ 2 ＋ 2} を、 かけ\R{算}{ざん}の しきに しよう。
\bkshikiran[3\zw]
\bkpans{\mbox{2 × 7 ＝ 14}\quad 14}

\bkproblem \mbox{3 × 6 ＝ 18} です。 \mbox{3 × 7} は いくつ？\par\vspace{2pt}
\noindent\hspace*{1\zw}\mbox{3 × 7 ＝ 18 ＋} \bkkakikomi\enspace ＝ \bkkakikomi
\bkpans{\mbox{3 × 7 ＝ 18 ＋ 3 ＝ 21}}

\clearpage
\bkproblemgroup{チャレンジ}{1}

\bkproblem 4\,cmの テープを 8\R{本}{ほん} つなぐと、 なん\,cm？
\bkshikiran
\bkpans{\mbox{4 × 8 ＝ 32}\quad 32\,cm}

\bkproblem 2こずつ 8さらと、 3こずつ 5さら。 おおいのは どっち？ なんこ おおい？\enspace{\small\color{bkOrange!58!black}ヒント：ちがいは ひき\R{算}{ざん}（1-5）}
\bkshikiran[3\zw]
\bkshikiran[3\zw]
\par\vspace{2pt}\noindent\hfill\bkkakikomi[28mm]\enspace の ほうが \bkkakikomi\enspace こ おおい。
\bkpans{\mbox{2 × 8 ＝ 16}、 \mbox{3 × 5 ＝ 15}\quad 2こずつ 8さらの ほうが 1こ おおい}

\begin{furikaeri}\bkqkrag
\begin{checklist}
\item 「1つ\R{分}{ぶん} × いくつ\R{分}{ぶん}」の しきが かける。\hfill{\footnotesize\sffamily\color{bkTheme}→\pageref{bk:1-12:matomari}ページ}
\item かけ\R{算}{ざん}を おはじきの \R{四角}{し|かく}で あらわせる。\hfill{\footnotesize\sffamily\color{bkTheme}→\pageref{bk:1-12:shikaku}ページ}
\item 2のだん〜5のだんの \R{九九}{く|く}を つかえる。\hfill{\footnotesize\sffamily\color{bkTheme}→\pageref{bk:1-12:dan}ページ}
\end{checklist}
チェックが つかなかったら、 →の ページに もどろう。
\end{furikaeri}

\bkfukidashi{sora}{おはじきの \mbox{\R{四角}{し|かく}}は、 1-13の \R{九九}{く|く}の \R{表}{ひょう}や、 1-21の \R{広}{ひろ}さに つながるよ。 みの まわりの 「〜ずつ」も さがして みよう！}
\clearpage
